\documentclass[10pt,a4paper]{amsart}
\usepackage{amsmath,amssymb,amsthm,graphicx}
% Retained minimal preamble for the Stacks Project chapter more-morphisms.tex (More on Morphisms).
% Verbatim excerpt of the chapter preamble: Xy-pic setup (native preamble lines 14-18) and the native
% theorem declarations with the mathematical macros (native preamble lines 154-217). Copied mechanically;
% nothing added, removed or redefined. The wrapper class and the amsmath/amssymb/amsthm/graphicx base
% packages are supplied by the assembler, so the class/externaldocument part of the native preamble is omitted.
% For commutative diagrams we use Xy-pic
\usepackage[all]{xy}

% We use 2cell for 2-commutative diagrams.
\xyoption{2cell}
% Theorem environments.
%
\theoremstyle{plain}
\newtheorem{theorem}[subsection]{Theorem}
\newtheorem{proposition}[subsection]{Proposition}
\newtheorem{lemma}[subsection]{Lemma}

\theoremstyle{definition}
\newtheorem{definition}[subsection]{Definition}
\newtheorem{example}[subsection]{Example}
\newtheorem{exercise}[subsection]{Exercise}
\newtheorem{situation}[subsection]{Situation}

\theoremstyle{remark}
\newtheorem{remark}[subsection]{Remark}
\newtheorem{remarks}[subsection]{Remarks}

\numberwithin{equation}{subsection}

% Macros
%
\def\lim{\mathop{\mathrm{lim}}\nolimits}
\def\colim{\mathop{\mathrm{colim}}\nolimits}
\def\Spec{\mathop{\mathrm{Spec}}}
\def\Hom{\mathop{\mathrm{Hom}}\nolimits}
\def\Ext{\mathop{\mathrm{Ext}}\nolimits}
\def\SheafHom{\mathop{\mathcal{H}\!\mathit{om}}\nolimits}
\def\SheafExt{\mathop{\mathcal{E}\!\mathit{xt}}\nolimits}
\def\Sch{\mathit{Sch}}
\def\Mor{\mathop{\mathrm{Mor}}\nolimits}
\def\Ob{\mathop{\mathrm{Ob}}\nolimits}
\def\Sh{\mathop{\mathit{Sh}}\nolimits}
\def\NL{\mathop{N\!L}\nolimits}
\def\CH{\mathop{\mathrm{CH}}\nolimits}
\def\proetale{{pro\text{-}\acute{e}tale}}
\def\etale{{\acute{e}tale}}
\def\QCoh{\mathit{QCoh}}
\def\Ker{\mathop{\mathrm{Ker}}}
\def\Im{\mathop{\mathrm{Im}}}
\def\Coker{\mathop{\mathrm{Coker}}}
\def\Coim{\mathop{\mathrm{Coim}}}

% Boxtimes
%
\DeclareMathSymbol{\boxtimes}{\mathbin}{AMSa}{"02}

%
% Macros for moduli stacks/spaces
%
\def\QCohstack{\mathcal{QC}\!\mathit{oh}}
\def\Cohstack{\mathcal{C}\!\mathit{oh}}
\def\Spacesstack{\mathcal{S}\!\mathit{paces}}
\def\Quotfunctor{\mathrm{Quot}}
\def\Hilbfunctor{\mathrm{Hilb}}
\def\Curvesstack{\mathcal{C}\!\mathit{urves}}
\def\Polarizedstack{\mathcal{P}\!\mathit{olarized}}
\def\Complexesstack{\mathcal{C}\!\mathit{omplexes}}
% \Pic is the operator that assigns to X its picard group, usage \Pic(X)
% \Picardstack_{X/B} denotes the Picard stack of X over B
% \Picardfunctor_{X/B} denotes the Picard functor of X over B
\def\Pic{\mathop{\mathrm{Pic}}\nolimits}
\def\Picardstack{\mathcal{P}\!\mathit{ic}}
\def\Picardfunctor{\mathrm{Pic}}
\def\Deformationcategory{\mathcal{D}\!\mathit{ef}}
\title{A morphism of affine schemes is a closed immersion exactly when it lifts against every injective ring map}
\author{The Stacks Project Authors}
\date{stacks-project:more-morphisms.tex:local-snapshot-20261011; GNU FDL 1.2 or later}
\begin{document}
\maketitle
\noindent\textit{Editorial note.} Statement span is the whole lemma environment lines 23495-23507, including the commutative square with the dotted lift; proof span is lines 23509-23544.\par\medskip
\noindent\textit{Editorial note.} The proof contains two xymatrix displays; the retained preamble provides the Xy-pic setup.\par\medskip
\noindent\textit{Editorial note.} The proof cites no label and no bibliography entry, so no external\_references are needed.\par\medskip
\noindent\textit{Editorial note.} The original text contains the typographical quirk '\textbackslash{}medskip\textbackslash{}noindent' between paragraphs and the phrase 'and we win'; both are retained verbatim.\par\medskip
\noindent\textit{Editorial note.} extra\_packages is empty on purpose: the retained preamble already loads xy with options and a second declaration clashes.\par\medskip
\noindent\textit{Allowed standard prerequisites (editorial).}\begin{itemize}
\item Explicit prerequisite boundary. The dictionary between closed immersions of affine schemes and surjective ring maps, pullback squares of spectra, and the Yoneda/anti-equivalence bookkeeping. All are standard; both directions of the criterion are retained.
\end{itemize}\medskip
\begin{lemma}
\label{lemma-affine-extensive-criterion}
A morphism $f : X \to Y$ of affine schemes is a closed immersion 
if and only if for every injective ring map $A \to B$ and commutative 
square
$$
\xymatrix{
\Spec(B) \ar[d] \ar[r] & X \ar[d]^f \\
\Spec(A) \ar[r] \ar@{..>}[ur] & Y
}
$$
there exists a lift $\Spec(A) \to X$ making the two triangles commute.
\end{lemma}

\begin{proof}
Let the morphism $f$ be given by the ring map $\phi : R \to S$.
Then $f$ is a closed immersion if and only if $\phi$ is surjective.

\medskip\noindent
First, we assume that $\phi$ is surjective.
Let $\psi : A \to B$ be an injective ring map, and suppose we are given 
a commutative diagram
$$
\xymatrix{
R \ar[r]^\alpha \ar[d]^\phi & A \ar[d]^\psi \\
S \ar[r]^\beta \ar@{..>}[ur] & B
}
$$
Then we define a lift $S \to A$ by $s \mapsto \alpha(r)$, where 
$r \in R$ is such that $\phi(r) = s$. 
This is well-defined because $\psi$ is injective and the square commutes.
Since taking the ring spectrum defines an anti-equivalence between 
commutative rings and affine schemes, the desired lifting property for 
$f$ holds.

\medskip\noindent
Next, we assume that $\phi$ has lifts against all injective ring 
maps $\psi: A \to B$.
Note that $\phi(R)$ is a subring of $S$, so we obtain a commutative 
square
$$
\xymatrix{
R \ar[r] \ar[d]^\phi  & \phi(R) \ar[d] \\
S \ar@{=}[r] \ar@{..>}[ur] & S
}
$$
in which a lift $S \to \phi(R)$ exists. 
Hence, the inclusion $\phi(R) \to S$ must be an isomorphism, which 
shows that $\phi$ is surjective, and we win.
\end{proof}

\end{document}
